% !atlas-meta
% id: upanishads
% title: Upanishads
% author: badarayana
% authorName: Various
% language: English
% sourceLanguage: Sanskrit
% translator: F. Max Müller
% translationYear: 1879
% source: https://en.wikisource.org/wiki/Sacred_Books_of_the_East/Volume_1
% license: Public Domain (translator F. Max Müller d. 1900)
% status: complete
% !end-atlas-meta
\documentclass[12pt]{article}
\usepackage[utf8]{inputenc}
\usepackage{ebgaramond}
\usepackage[margin=1.2in]{geometry}
\usepackage{parskip}
\newcommand{\work}[1]{\begin{center}\LARGE\scshape #1\end{center}}
\newcommand{\attrib}[2]{\begin{center}\large #1\\[2pt]\small\itshape trans.\ #2\end{center}\vspace{1em}}
\newcommand{\para}[1][]{\par\noindent\ifx&#1&\else\textsuperscript{\footnotesize #1}~\fi}
\begin{document}
% !body-start
\work{The Upanishads}
\attrib{Anonymous (Vedic sages)}{F. Max M\"uller (public domain)}

\section*{Isha-Upanishad (V\^agasaneyi-Samhit\^a)}

\para[1] All this, whatsoever moves on earth, is to be hidden in the Lord (the Self). When thou hast surrendered all this, then thou mayest enjoy. Do not covet the wealth of any man!

\para[2] Though a man may wish to live a hundred years, performing works, it will be thus with him; but not in any other way: work will thus not cling to a man.

\para[3] There are the worlds of the Asuras covered with blind darkness. Those who have destroyed their self (who perform works, without having arrived at a knowledge of the true Self), go after death to those worlds.

\para[4] That one (the Self), though never stirring, is swifter than thought. The Devas (senses) never reached it, it walked before them. Though standing still, it overtakes the others who are running. Mâtarisvan (the wind, the moving spirit) bestows powers on it.

\para[5] It stirs and it stirs not; it is far, and likewise near. It is inside of all this, and it is outside of all this.

\para[6] And he who beholds all beings in the Self, and the Self in all beings, he never turns away from it.

\para[7] When to a man who understands, the Self has become all things, what sorrow, what trouble can there be to him who once beheld that unity?

\para[8] He (the Self) encircled all, bright, incorporeal, scatheless, without muscles, pure, untouched by evil; a seer, wise, omnipresent, self-existent, he disposed all things rightly for eternal years.

\para[9] All who worship what is not real knowledge (good works), enter into blind darkness: those who delight in real knowledge, enter, as it were, into greater darkness.

\para[10] One thing, they say, is obtained from real knowledge; another, they say, from what is not knowledge. Thus we have heard from the wise who taught us this.

\para[11] He who knows at the same time both knowledge and not-knowledge, overcomes death through not-knowledge, and obtains immortality through knowledge.

\para[12] All who worship what is not the true cause, enter into blind darkness: those who delight in the true cause, enter, as it were, into greater darkness.

\para[13] One thing, they say, is obtained from (knowledge of) the cause; another, they say, from (knowledge of) what is not the cause. Thus we have heard from the wise who taught us this.

\para[14] He who knows at the same time both the cause and the destruction (the perishable body), overcomes death by destruction (the perishable body), and obtains immortality through (knowledge of) the true cause.

\para[15] The door of the True is covered with a golden disk. Open that, O Pûshan, that we may see the nature of the True.

\para[16] O Pûshan, only seer, Yama (judge), Sûrya (sun), son of Pragâpati, spread thy rays and gather them! The light which is thy fairest form, I see it. I am what He is (viz. the person in the sun).

\para[17] Breath to air, and to the immortal! Then this my body ends in ashes. Om! Mind, remember! Remember thy deeds! Mind, remember! Remember thy deeds

\para[18] Agni, lead us on to wealth (beatitude) by a good path, thou, O God, who knowest all things! Keep far from us crooked evil, and we shall offer thee the fullest praise!

\section*{Kena-Upanishad (Talavak\^ara)}

\section*{First Khanda}

\para[1] The Pupil asks: 'At whose wish does the mind sent forth proceed on its errand? At whose command does the first breath go forth? At whose wish do we utter this speech? What god directs the eye, or the ear?'

\para[2] The Teacher replies: It is the ear of the ear, the mind of the mind, the speech of speech, the breath of breath, and the eye of the eye. When freed (from the senses) the wise, on departing from this world, become immortal.

\para[3] 'The eye does not go thither, nor speech, nor mind. We do not know, we do not understand, how any one can teach it.

\para[4] 'It is different from the known, it is also above the unknown, thus we have heard from those of old, who taught us this.

\para[5] 'That which is not expressed by speech and by which speech is expressed, that alone know as Brahman, not that which people here adore.

\para[6] 'That which does not think by mind, and by which, they say, mind is thought, that alone know as Brahman, not that which people here adore.

\para[7] 'That which does not see by the eye, and by which one sees (the work of) the eyes, that alone know as Brahman, not that which people here adore.

\para[8] 'That which does not hear by the ear, and by which the ear is heard, that alone know as Brahman, not that which people here adore.

\para[9] 'That which does not breathe by breath, and by which breath is drawn, that alone know as Brahman, not that which people here adore.'

\section*{Second Khanda}

\para[1] The Teacher says: 'If thou thinkest I know it well; then thou knowest surely but little, what is that form of Brahman known, it may be, to thee?'

\para[2] The Pupil says: 'I do not think I know it well, nor do I know that I do not know it. He among us who knows this, he knows it, nor does he know that he does not know it.

\para[3] 'He by whom it (Brahman) is not thought, by him it is thought; he by whom it is thought, knows it not. It is not understood by those who understand it, it is understood by those who do not understand it.

\para[4] 'It is thought to be known (as if) by awakening, and (then) we obtain immortality indeed. By the Self we obtain strength, by knowledge we obtain immortality.

\para[5] 'If a man know this here, that is the true (end of life); if he does not know this here, then there is great destruction (new births). The wise who have thought on all things (and recognised the Self in them) become immortal, when they have departed from this world.'

\section*{Third Khanda}

\para[1] Brahman obtained the victory for the Devas. The Devas became elated by the victory of Brahman, and they thought, this victory is ours only, this greatness is ours only.

\para[2] Brahman perceived this and appeared to them. But they did not know it, and said: 'What sprite (yaksha or yakshya) is this?'

\para[3] They said to Agni (fire): 'O Gâtavedas, find out what sprite this is.' 'Yes,' he said.

\para[4] He ran toward it, and Brahman said to him: 'Who are you?' He replied: 'I am Agni, I am Gâtavedas.'

\para[5] Brahman said: 'What power is in you?' Agni replied: 'I could burn all whatever there is on earth.'

\para[6] Brahman put a straw before him, saying: 'Burn this.' He went towards it with all his might, but he could not burn it. Then he returned thence and said: 'I could not find out what sprite this is.'

\para[7] Then they said to Vâyu (air): 'O Vâyu, find out what sprite this is.' 'Yes,' he said.

\para[8] He ran toward it, and Brahman said to him: 'Who are you?' He replied: 'I am Vâyu, I am Mâtarisvan.'

\para[9] Brahman said: 'What power is in you?' Vâyu replied: 'I could take up all whatever there is on earth.'

\para[10] Brahman put a straw before him, saying: 'Take it up.' He went towards it with all his might, but he could not take it up. Then he returned thence and said: 'I could not find out what sprite this is.'

\para[11] Then they said to Indra: 'O Maghavan, find out what sprite this is.' He went towards it, but it disappeared from before him.

\para[12] Then in the same space (ether) he came towards a woman, highly adorned: it was Umâ, the daughter of Himavat. He said to her: 'Who is that sprite?'

\section*{Fourth Khanda}

\para[1] She replied: 'It is Brahman. It is through the victory of Brahman that you have thus become great.' After that he knew that it was Brahman.

\para[2] Therefore these Devas, viz. Agni, Vâyu, and Indra, are, as it were, above the other gods, for they touched it (the Brahman) nearest.

\para[3] And therefore Indra is, as it were, above the other gods, for he touched it nearest, he first knew it.

\para[4] This is the teaching of Brahman, with regard to the gods (mythological): It is that which now flashes forth in the lightning, and now vanishes again.

\para[5] And this is the teaching of Brahman, with regard to the body (psychological): It is that which seems to move as mind, and by it imagination remembers again and again.

\para[6] That Brahman is called Tadvana, by the name of Tadvana it is to be meditated on. All beings have a desire for him who knows this.

\para[7] The Teacher: 'As you have asked me to tell you the Upanishad, the Upanishad has now been told you. We have told you the Brâhmî Upanishad.

\para[8] 'The feet on which that Upanishad stands are penance, restraint, sacrifice; the Vedas are all its limbs, the True is its abode.

\para[9] 'He who knows this Upanishad, and has shaken off all evil, stands in the endless, unconquerable world of heaven, yea, in the world of heaven.'

\section*{Katha-Upanishad}

\section*{FIRST ADHYÂYA}

\section*{First Vallî}

\para[1] Vâgasravasa, desirous (of heavenly rewards), surrendered (at a sacrifice) all that he possessed. He had a son of the name of Nakiketas.

\para[2] When the (promised) presents were being given (to the priests), faith entered into the heart of Nakiketas, who was still a boy, and he thought:

\para[3] "Unblessed, surely, are the worlds to which a man goes by giving (as his promised present at a sacrifice) cows which have drunk water, eaten hay, given their milk, and are barren."

\para[4] He (knowing that his father had promised to give up all that he possessed, and therefore his son also) said to his father: "Dear father, to whom wilt thou give me?" He said it a second and a third time. Then the father replied (angrily): "I shall give thee unto Death." (The father, having once said so, though in haste, had to be true to his word and to sacrifice his son.)

\para[5] The son said: "I go as the first, at the head of many (who have still to die); I go in the midst of many (who are now dying). What will be the work of Yama (the ruler of the departed) which to-day he has to do unto me?

\para[6] Look back how it was with those who came before, look forward how it will be with those who come hereafter. A mortal ripens like corn, like corn he springs up again." (Nakiketas enters into the abode of Yama Vaivasvata, and there is no one to receive him. Thereupon one of the attendants of Yama is supposed to say:)

\para[7] "Fire enters into the houses, when a Brâhmana enters as a guest. That fire is quenched by this peace-offering; --- bring water, O Vaivasvata!

\para[8] A Brâhmana that dwells in the house of a foolish man without receiving food to eat, destroys his hopes and expectations, his possessions, his righteousness, his sacred and his good deeds, and all his sons and cattle." (Yama, returning to his house after an absence of three nights, during which time Nakiketas had received no hospitality from him, says:)

\para[9] "O Brâhmana, as thou, a venerable guest, hast dwelt in my house three nights without eating, therefore choose now three boons. Hail to thee! and welfare to me!"

\para[10] Nakiketas said: "O Death, as the first of the three boons I choose that Gautama, my father, be pacified, kind, and free from anger towards me; and that he may know me and greet me, when I shall have been dismissed by thee."

\para[11] Yama said: "Through my favour Auddâlaki Âruni, thy father, will know thee, and be again towards thee as he was before. He shall sleep peacefully through the night, and free from anger, after having seen thee freed from the mouth of death."

\para[12] Nakiketas said: "In the heaven-world there is no fear; thou art not there, O Death, and no one is afraid on account of old age. Leaving behind both hunger and thirst, and out of the reach of sorrow, all rejoice in the world of heaven.

\para[13] Thou knowest, O Death, the fire-sacrificef which leads us to heaven; tell it to me, for I am full of faith. Those who live in the heaven-world reach immortality,- this I ask as my second boon."

\para[14] Yama said: "I tell it thee, learn it from me, and when thou understandest that fire-sacrifice which leads to heaven, know, O Nakiketas, that it is the attainment of the endless worlds, and their firm support, hidden in darkness."

\para[15] Yama then told him that fire-sacrifice, the beginning of all the worlds, and what bricks are required for the altar, and how many, and how they are to be placed. And Nakiketas repeated all as it had been told to him. Then Mrityu, being pleased with him, said again:

\para[16] The generous, being satisfied, said to him: "I give thee now another boon; that fire-sacrifice shall be named after thee, take also this many-coloured chain.

\para[17] He who has three times performed this Nâkiketa-rite and has been united with the three (father, mother, and teacher), and has performed the three duties (study, sacrifice, almsgiving) overcomes birth and death. When he has learnt and understood this fire, which knows (or makes us know) all that is born of Brahman, which is venerable and divine, then he obtains everlasting peace.

\para[18] He who knows the three Nâkiketa fires, and knowing the three, piles up the Nâkiketa sacrifice, he, having first, thrown off the chains of death, rejoices in the world of heaven, beyond the reach of grief.

\para[19] This, O Nakiketa, is thy fire which leads to heaven, and which thou hast chosen as thy second boon. That fire all men will proclaim. Choose now, O Nakiketa, thy third boon."

\para[20] Nakiketa said: "There is that doubt, when a man is dead, --- some saying, he is; others, he is not. This I should like to know, taught by thee; this is the third of my boons."

\para[21] Death said: "On this point even the gods have doubted formerly; it is not easy to understand. That subject is subtle. Choose another boon, O Nakiketas, do not press me, and let me off that boon."

\para[22] Nakiketas said: "On this point even the gods have doubted indeed, and thou, Death, hast declared it to be not easy to understand, and another teacher like thee is not to be found: --- surely no other boon is like unto this."

\para[23] Death said: "Choose sons and grandsons who shall live a hundred years, herds of cattle, elephants, gold, and horses. Choose the wide abode of the earth, and live thyself as many harvests as thou desirest.

\para[24] If you can think of any boon equal to that, choose wealth, and long life. Be (king), Nakiketas, on the wide earth. I make thee the enjoyer of all desires.

\para[25] Whatever desires are difficult to attain among mortals, ask for them according to thy wish; --- these fair maidens with their chariots and musical instruments, --- such are indeed not to be obtained by men, --- be waited on by them whom I give to thee, but do not ask me about dying."

\para[26] Nakiketas said: "These things last till tomorrow, O Death, for they wear out this vigour of all the senses. Even the whole of life is short. Keep thou thy horses, keep dance and song for thyself.

\para[27] No man can be made happy by wealth. Shall we possess wealth, when we see thee? Shall we live, as long as thou rulest? Only that boon (which I have chosen) is to be chosen by me.

\para[28] What mortal, slowly decaying here below, and knowing, after having approached them, the freedom from decay enjoyed by the immortals, would delight in a long life, after he has pondered on the pleasures which arise from beauty and love?

\para[29] No, that on which there is this doubt, O Death, tell us what there is in that great Hereafter. Nakiketas does not choose another boon but that which enters into the hidden world."

\section*{SECOND Vallî}

\para[1] Death said: "The good is one thing, the pleasant another; these two, having different objects, chain a man. It is well with him who clings to the good; he who chooses the pleasant, misses his end.

\para[2] The good and the pleasant approach man: the wise goes round about them and distinguishes them. Yea, the wise prefers the good to the pleasant, but the fool chooses the pleasant through greed and avarice.

\para[3] Thou, O Nakiketas, after pondering all pleasures that are or seem delightful, hast dismissed them all. Thou hast not gone into the road that leadeth to wealth, in which many men perish.

\para[4] Wide apart and leading to different points are these two, ignorance, and what is known as wisdom. I believe Nakiketas to be one who desires knowledge, for even many pleasures did not tear thee away.

\para[5] Fools dwelling in darkness, wise in their own conceit, and puffed up with vain knowledge, go round and round, staggering to and fro, like blind men led by the blind.

\para[6] The Hereafter never rises before the eyes of the careless child, deluded by the delusion of wealth. "This is the world," he thinks, "there is no other;" --- thus he falls again and again under my sway.

\para[7] He (the Self) of whom many are not even able to hear, whom many, even when they hear of him, do not comprehend; wonderful is a man, when found, who is able to teach him (the Self); wonderful is he who comprehends him, when taught by an able teacher.

\para[8] That (Self), when taught by an inferior man, is not easy to be known, even though often thought upon; unless it be taught by another, there is no way to it, for it is inconceivably smaller than what is small.

\para[9] That doctrine is not to be obtained by argument, but when it is declared by another, then, O dearest, it is easy to understand. Thou hast obtained it now; thou art truly a man of true resolve. May we have always an inquirer like thee!"

\para[10] Nakiketas said: "I know that what is called a treasure is transient, for that eternal is not obtained by things which are not eternal. Hence the Nâkiketa fire(sacrifice) has been laid by me (first); then by means of transient things, I have obtained what is not transient (the teaching of Yama)."

\para[11] Yama said: "Though thou hadst seen the fulfilment of all desires, the foundation of the world, the endless rewards of good deeds, the shore where there is no fear, that which is magnified by praise, the wide abode, the rest, yet being wise thou hast with firm resolve dismissed it all.

\para[12] The wise who, by means of meditation on his Self, recognises the Ancient, who is difficult to be seen, who has entered into the dark, who is hidden in the cave, who dwells in the abyss, as God, he indeed leaves joy and sorrow far behind."

\para[13] A mortal who has heard this and embraced it, who has separated from it all qualities, and has thus reached the subtle Being, rejoices, because he has obtained what is a cause for rejoicing. The house (of Brahman) is open, I believe, O Nakiketas."

\para[14] Nakiketas said: "That which thou seest as neither this nor that, as neither effect nor cause, as neither past nor future, tell me that."

\para[15] Yama said: "That word (or place) which all - the Vedas record, which all penances proclaim, which men desire when they live as religious students, that word I tell thee briefly, it is Om.

\para[16] That (imperishable) syllable means Brahman, that syllable means the highest (Brahman); he who knows that syllable, whatever he desires,is his.

\para[17] This is the best support, this is the highest support; he who knows that support is magnified in the world of Brahmâ.

\para[18] The knowing (Self) is not born, it dies not; it sprang from nothing, nothing sprang from it. The Ancient is unborn, eternal, everlasting; he is not killed, though the body is killed.

\para[19] If the killer thinks that he kills, if the killed thinks that he is killed, they do not understand; for this one does not kill, nor is that one killed.

\para[20] The Self, smaller than small, greater than great, is hidden in the heart of that creature. A man who is free from desires and free from grief, sees the majesty of the Self by the grace of the Creator.

\para[21] Though sitting still, he walks far; though lying down, he goes everywhere. Who, save myself, is able to know that God who rejoices and rejoices not?

\para[22] The wise who knows the Self as bodiless within the bodies, as unchanging among changing things, as great and omnipresent, does never grieve.

\para[23] That Self cannot be gained by the Veda, nor by understanding, nor by much learning. He whom the Self chooses, by him the Self can be can gained. The Self chooses him (his body) as his own.

\para[24] But he who has not first turned away from his wickedness, who is not tranquil, and subdued, or whose mind is not at rest, he can never obtain the Self (even) by knowledge.

\para[25] Who then knows where He is, He to whom the Brahmans and Kshatriyas are (as it were) but food, and death itself a condiment?"

\section*{THIRD Vallî}

\para[1] "There are the two, drinking their reward in the world of their own works, entered into the cave (of the heart), dwelling on the highest summit (the ether in the heart). Those who know Brahman call them shade and light; likewise, those householders who perform the Trinâkiketa sacrifice.

\para[2] May we be able to master that Nâkiketa rite which is a bridge for sacrificers; also that which is the highest, imperishable Brahman for those who wish to cross over to the fearless shore.

\para[3] Know the Self to be sitting in the chariot, the body to be the chariot, the intellect (buddhi) the charioteer, and the mind the reins.

\para[4] The senses they call the horses, the objects of the senses their roads. When he (the Highest Self) is in union with the body, the senses, and the mind, then wise people call him the Enjoyer.

\para[5] He who has no understanding and whose mind (the reins) is never firmly held, his senses (horses) are unmanageable, like vicious horses of a charioteer.

\para[6] But he who has understanding and whose mind is always firmly held, his senses are under control, like good horses of a charioteer.

\para[7] He who has no understanding, who is unmindful and always impure, never reaches that place, but enters into the round of births.

\para[8] But he who has understanding, who is mindful and always pure, reaches indeed that place; from whence he is not born again.

\para[9] But he who has understanding for his charioteer, and who holds the reins of the mind, he reaches the end of his journey, and that is the highest place of Vishnu.

\para[10] Beyond the senses there are the objects, beyond the objects there is the mind, beyond the mind there is the intellect, the Great Self is beyond the intellect.

\para[11] Beyond the Great there is the Undeveloped, beyond the Undeveloped there is the Person (purusha). Beyond the Person there is nothing --- this is the goal, the highest road.

\para[12] That Self is hidden in all beings and does not shine forth, but it is seen by subtle seers through their sharp and subtle intellect.

\para[13] A wise man should keep down speech and mind; he should keep them within the Self which is knowledge; he should keep knowledge within the Self whichs is the Great; and he should keep that (the Great) within the Self which is the Quiet.

\para[14] Rise, awake! having obtained your boons, understand them! The sharp edge of a razor is difficult to pass over; thus the wise say the path (to the Self) is hard.

\para[18] He who has perceived that which is without sound, without touch, without form, without decay, without taste, eternal, without smell, without beginning, without end, beyond the Great, and unchangeable, is freed from the jaws of death.

\para[19] A wise man who has repeated or heard the ancient story of Nakiketas told by Death, is magnified in the world of Brahman.

\para[20] And he who repeats this greatest mystery in an assembly of Brâhmans, or full of devotion at the time of the Srâddha sacrifice, obtains thereby infinite rewards."

\section*{SECOND ADHYÂYA}

\section*{FOURTH Vallî}

\para[1] Death said: "The Self-existent pierced the openings (of the senses) so that they turn forward; therefore man looks forward, not backward into himself. Some wise man, however, with his eyes closed and wishing for immortality, saw the Self behind.

\para[2] Children follow after outward pleasures, and fall into the snare of wide-spread death. Wise men only, knowing the nature of what is immortal, do not look for anything stable here among things unstable.

\para[3] That by which we know form, taste, smell, sounds, and loving touches, by that also we know what exists besides. This is that (which thou hast asked for).

\para[4] The wise, when he knows that that by which he perceives all objects in sleep or in waking is the great omnipresent Self, grieves no more.

\para[5] He who knows this living soul, which eats honey (perceives objects) as being the Self, always near, the Lord of the past and the future, henceforward fears no more. This is that.

\para[6] He who (knows) him who was born first from the brooding heat (for he was born before the water), who, entering into the heart, abides therein, and was perceived from the elements. This is that.

\para[7] (He who knows) Aditi also, who is one with all deities, who arises with Prâna (breath or Hiranyagarbha), who, entering into the heart, abides therein, and was born from the elements. This is that.

\para[8] There is Agni (fire), the all-seeing, hidden in the two fire-sticks, well-guarded like a child (in the womb) by the mother, day after day to be adored by men when they awake and bring oblations. This is that.

\para[9] And that whence the sun rises, and whither it goes to set, there all the Devas are contained, and no one goes beyond. This is that.

\para[10] What is here (visible in the world), the same is there (invisible in Brahman); and what is there, the same is here. He who sees any difference here (between Brahman and the world), goes from death to death.

\para[11] Even by the mind this (Brahman) is to be obtained, and then there is no difference whatsoever. He goes from death to death who sees any difference here.

\para[12] The person (purusha), of the size of a thumb, stands in the middle of the Self (body ?), as lord of the past and the future, and henceforward fears no more. This is that.

\para[13] That person, of the size of a thumb, is like a light without smoke, lord of the past and the future, he is the same to-day and to-morrow. This is that.

\para[14] As rain-water that has fallen on a mountain-ridge runs down the rocks on all sides, thus does he, who sees a difference between qualities, run after them on all sides.

\para[15] As pure water poured into pure water remains the same, thus, O Gautama, is the Self of a thinker who knows."

\section*{FIFTH Vallî}

\para[1] "There is a town with eleven gates belonging to the Unborn (Brahman), whose thoughts are never crooked. He who approaches it, grieves no more, and liberated (from all bonds of ignorance) becomes free. This is that.

\para[2] He (Brahman) is the swan (sun), dwelling in the bright heaven; he is the Vasu (air), dwelling in the sky; he is the sacrificer (fire), dwelling on the hearth; he is the guest (Soma), dwelling in the sacrificial jar; he dwells in men, in gods (vara), in the sacrifice (rita), in heaven; he is born in the water, on earth, in the sacrifice (rita), on the mountains; he is the True and the Great.

\para[3] He (Brahman) it is who sends up the breath (prâna), and who throws back the breath (apâna). All the Devas (senses) worship him, the adorable (or the dwarf), who sits in the centre.

\para[4] When that incorporated (Brahman), who dwells in the body, is torn away and freed from the body, what remains then? This is that.

\para[5] No mortal lives by the breath that goes up and by the breath that goes down. We live by another, in whom these two repose.

\para[6] Well then, O Gautama, I shall tell thee this mystery, the old Brahman, and what happens to the Self, after reaching death.

\para[7] Some enter the womb in order to have a body, as organic beings, others go into inorganic matter, according to their work and according to their knowledge.

\para[8] He, the highest Person, who is awake in us while we are asleep, shaping one lovely sight after another, that indeed is the Bright, that is Brahman, that alone is called the Immortal. All worlds are contained in it, and no one goes beyond. This is that.

\para[9] As the one fire, after it has entered the world, though one, becomes different according to whatever it burns, thus the one Self within all things becomes different, according to whatever it enters, and exists also without.

\para[10] As the one air, after it has entered the world, though one, becomes different according to whatever it enters, thus the one Self within all things becomes different, according to whatever it enters, and exists also without.

\para[11] As the sun, the eye of the whole world, is not contaminated by the external impurities seen by the eyes, thus the one Self within all things is never contaminated by the misery of the world, being himself without.

\para[12] There is one ruler, the Self within all things, who makes the one form manifold. The wise who perceive him within their Self, to them belongs eternal happiness, not to others.

\para[13] There is one eternal thinker, thinking non-eternal thoughts, who, though one, fulfils the desires of many. The wise who perceive him within their Self, to them belongs eternal peace, not to others.

\para[14] They perceive that highest indescribable pleasure, saying, This is that. How then can I understand it? Has it its own light, or does it reflect light?

\para[15] The sun does not shine there, nor the moon and the stars, nor these lightnings, and much less this fire. When he shines, everything shines after him; by his light all this is lighted."

\section*{SIXTH Vallî}

\para[1] "There is that ancient tree, whose roots grow upward and whose branches grow downward; --- that indeed is called the Bright that is called Brahman, that alone is called the Immortal. All worlds are contained in it, and no one goes beyond. This is that.

\para[2] Whatever there is, the whole world, when gone forth (from the Brahman), trembles in its breath. That Brahman is a great terror, like a drawn sword. Those who know it become immortal.

\para[3] From terror of Brahman fire burns, from terror the sun burns, from terror Indra and Vâyu, and Death, as the fifth, run away.

\para[4] If a man could not understand it before the falling asunder of his body, then he has to take body again in the worlds of creation.

\para[5] As in a mirror, so (Brahman may be seen clearly) here in this body; as in a dream, in the world of the Fathers; as in the water, he is seen about in the world of the Gandharvas; as in light and shade in the world of Brahmâ.

\para[6] Having understood that the senses are distinct (from the Âtman), and that their rising and setting (their waking and sleeping) belongs to them in their distinct existence (and not to the Âtman), a wise man grieves no more.

\para[7] Beyond the senses is the mind, beyond the mind is the highest (created) Being, higher than that Being is the Great Self, higher than the Great, the highest Undeveloped.

\para[8] Beyond the Undeveloped is the Person, the all-pervading and entirely imperceptible. Every creature that knows him is liberated, and obtains immortality.

\para[9] His form is not to be seen, no one beholds him with the eye. He is imagined by the heart, by wisdom, by the mind. Those who know this, are immortal.

\para[10] When the five instruments of knowledge stand still together with the mind, and when the intellect does not move, that is called the highest state.

\para[11] This, the firm holding back of the senses, is what is called Yoga. He must be free from thoughtlessness then, for Yoga comes and goes.

\para[12] He (the Self) cannot be reached by speech, by mind, or by the eye. How can it be apprehended except by him who says: "He is?"

\para[13] By the words "He is," is he to be apprehended, and by (admitting) the reality of both (the invisible Brahman and the visible world, as coming from Brahman). When he has been apprehended by the words "He is," then his reality reveals itself.

\para[14] When all desires that dwell in his heart cease, then the mortal becomes immortal, and obtains Brahman.

\para[15] Whean all the ties of the heart are severed here on earth, then the mortal becomes immortal --- here ends the teaching.

\para[16] There are a hundred and one arteries of the heart, one of them penetrates the crown of the head. Moving upwards by it, a man (at his death) reaches the Immortal; the other arteries serve for departing in different directions.

\para[17] The Person not larger than a thumb, the inner Self, is always settled in the heart of men. Let a man draw that Self forth from his body with steadiness, as one draws the pith from a reed. Let him know that Self as the Bright, as the Immortal; yes, as the Bright, as the Immortal.

\para[18] Having received this knowledge taught by Death and the whole rule of Yoga (meditation), Nakiketa became free from passion and death, and obtained Brahman. Thus it will be with another also who knows thus what relates to the Self.

\para[19] May He protect us both! May He enjoy us both! May we acquire strength together! May our knowledge become bright! May we never quarrel! Om! Peace! peace! peace! Harih, Om!"

\section*{Prasna-Upanishad}

\section*{FIRST QUESTION}

\para Adoration to the Highest Self! Harih, Om!

\para[1] Sukesas Bhâradvâga, and Saivya Satyakâma, and Sauryâyanin Gârgya, and Kausalya Âsvalâyana, and Bhârgava Vaidarbhi, and Kabandhin Kâtyâyana, these were devoted to Brahman, firm in Brahman, seeking for the Highest Brahman. They thought that the venerable Pippalâda could tell them all that, and they therefore took fuel in their hands (like pupils), and approached him.

\para[2] That Rishi said to them: "Stay here a year longer, with penance, abstinence, and faith; then you may ask questions according to your pleasure, and if we know them, we shall tell you all."

\para[3] Then Kabandhin Kâtyâyana approached him and asked: "Sir, from whence may these creatures be born?"

\para[4] He replied: "Pragâpati (the lord of creatures) was desirous of creatures (pragâh). He performed penance, and having performed penance, he produces a pair, matter (rayi) and spirit (prâna), thinking that they together should produce creatures for him in many ways.

\para[5] The sun is spirit, matter is the moon. All this, what has body and what has no body, is matter, and therefore body indeed is matter.

\para[6] Now Âditya, the sun, when he rises, goes toward the East, and thus receives the Eastern spirits into his rays. And when he illuminates the South, the West, the North, the Zenith, the Nadir, the intermediate quarters, and everything, he thus receives all spirits into his rays.

\para[7] Thus he rises, as Vaisvânara, (belonging to all men,) assuming all forms, as spirit, as fire. This has been said in the following verse:

\para[8] (They knew) him who assumes all forms, the golden, who knows all things, who ascends highest, alone in his splendour, and warms us; the thousand-rayed, who abides in a hundred places, the spirit of all creatures, the Sun, rises.

\para[9] The year indeed is Pragâpati, and there are two paths thereof, the Southern and the Northern. Now those who here believe in sacrifices and pious gifts as work done, gain the moon only as their (future) world, and return again. Therefore the Rishis who desire offspring, go to the South, and that path of the Fathers is matter (rayi).

\para[10] But those who have sought the Self by penance, abstinence, faith, and knowledge, gain by the Northern path Âditya, the sun. This is the home of the spirits, the immortal, free from danger, the highest. From thence they do not return, for it is the end. Thus says the Sloka:

\para[11] Some call him the father with five feet (the five seasons), and with twelve shapes (the twelve months), the giver of rain in the highest half of heaven; others again say that the sage is placed in the lower half, in the chariot with seven wheels and six spokes.

\para[12] The month is Pragâpati; its dark half is matter, its bright half spirit. Therefore some Rishis perform sacrifice in the bright half, others in the other half.

\para[13] Day and Night are Pragâpati; its day is spirit, its night matter. Those who unite in love by day waste their spirit, but to unite in love by night is right.

\para[14] Food is Pragâpati. Hence proceeds seed, and from it these creatures are born.

\para[15] Those therefore who observe this rule of Pragâpati (as laid down in §13), produce a pair, and to them belongs this Brahma-world here. But those in whom dwell penance, abstinence, and truth,

\para[16] To them belongs that pure Brahma-world, to them, namely, in whom there is nothing crooked, nothing false, and no guile."

\section*{SECOND QUESTION}

\para[1] Then Bhârgava Vaidarbhi asked him: "Sir, How many gods keep what has thus been created, how many manifest this, and who is the best of them?"

\para[2] He replied: "The ether is that god, the wind, fire, water, earth, speech, mind, eye, and ear. These, when they have manifested (their power), contend and say: We (each of us) support this body and keep it.

\para[3] Then Prâna (breath, spirit, life), as the best, said to them: Be not deceived, I alone, dividing myself fivefold, support this body and keep it.

\para[4] They were incredulous; so he, from pride, did as if he were going out from above. Thereupon, as he went out, all the others went out, and as he returned, all the others returned. As bees go out when their queen goes out, and return when she returns, thus (did) speech, mind, eye, and ear; and, being satisfied, they praise Prâna, saying:

\para[5] He is Agni (fire), he shines as Sûrya (sun), he is Parganya (rain), the powerful (Indra), he is Vâyu (wind), he is the earth, he is matter, he is God --- he is what is and what is not, and what is immortal.

\para[6] As spokes in the nave of a wheel, everything is fixed in Prâna, the verses of the Rig-veda, Yagur-veda, Sâma-veda, the sacrifice, the Kshatriyas, and the Brâhmans.

\para[7] As Pragâpati (lord of creatures) thou movest about in the womb, thou indeed art born again. To thee, the Prâna, these creatures bring offerings, to thee who dwellest with the other prânas (the organs of sense).

\para[8] Thou art the best carrier for the Gods, thou art the first offering to the Fathers. Thou art the true work of the Rishis, of the Atharvângiras.

\para[9] O Prâna, thou art Indra by thy light, thou art Rudra, as a protector; thou movest in the sky, thou art the sun, the lord of lights.

\para[10] When thou showerest down rain, then, O Prâna, these creatures of thine are delighted, hoping that there will be food, as much as they desire.

\para[11] Thou art a Vrâtya, O Prâna, the only Rishi, the consumer of everything, the good lord. We are the givers of what thou hast to consume, thou, O Mâtarisva, art our father.

\para[12] Make propitious that body of thine which dwells in speech, in the ear, in the eye, and which pervades the mind; do not go away!

\para[13] All this is in the power of Prâna, whatever exists in the three heavens. Protect us like a mother her sons, and give us happiness and wisdom."

\section*{THIRD QUESTION}

\para[1] Then Kausalya Âsvalâyana asked: "Sir, whence is that Prâna (spirit) born? How does it come into this body? And how does it abide, after it has divided itself? How does it go out? How does it support what is without, and how what is within?"

\para[2] He replied: "You ask questions more difficult, but you are very fond of Brahman, therefore I shall tell it you.

\para[3] This Prâna (spirit) is born of the Self. Like the shadow thrown on a man, this (the prâna) is spread out over it (the Brahman). By the work of the mind does it come into this body.

\para[4] As a king commands officials, saying to them: Rule these villages or those, so does that Prâna (spirit) dispose the other prânas, each for their separate work.

\para[5] The Apâna (the down-breathing) in the organs of excretion and generation; the Prâna himself dwells in eye and ear, passing through mouth and nose. In the middle is the Samâna (the on-breathing); it carries what has been sacrificed as food equally (over the body), and the seven lights proceed from it.

\para[6] The Self is in the heart. There are the 101 arteries, and in each of them there are a hundred (smaller veins), and for each of these branches there are 72,000. In these the Vyâna (the back-breathing) moves.

\para[7] Through one of them, the Udâna (the out-breathing) leads (us) upwards to the good world by good work, to the bad world by bad work, to the world of men by both.

\para[8] The sun rises as the external Prâna, for it assists the Prâna in the eye. The deity that exists in the earth, is there in support of man's Apâna (down-breathing). The ether between (sun and earth) is the Samâna (on-breathing), the air is Vyâna (back-breathing).

\para[9] Light is the Udâna (out-breathing), and therefore he whose light has gone out comes to a new birth with his senses absorbed in the mind.

\para[10] Whatever his thought (at the time of death) with that he goes back to Prâna, and the Prâna, united with light, together with the self (the gîvâtmâ) leads on to the world, as deserved.

\para[11] He who, thus knowing, knows Prâna, his offspring does not perish, and he becomes immortal. Thus says the Sloka:

\para[12] He who has known the origin, the entry, the place, the fivefold distribution, and the internal state of the Prâna, obtains immortality, yes, obtains immortality."

\section*{FOURTH QUESTION}

\para[1] Then Sauryâyanin Gârgya asked: "Sir, What are they that sleep in this man, and what are they that are awake in him? What power (deva) is it that sees dreams? Whose is the happiness? On what do all these depend?"

\para[2] He replied: "O Gârgya, As all the rays of the sun, when it sets, are gathered up in that disc of light, and as they, when the sun rises again and again, come forth, so is all this (all the senses) gathered up in the highest faculty (deva) the mind. Therefore at that time that man does not hear, see, smell, taste, touch, he does not speak, he does not take, does not enjoy, does not evacuate, does not move about. He sleeps, that is what people say.

\para[3] The fires of the prânas are, as it were, awake in that town (the body). The Apâna is the Gârhapatya fire, the Vyâna the Anvâhâryapakana fire; and because it is taken out of the Gârhapatya fire, which is fire for taking out, therefore the Prâna is the Âhavanlya fire.

\para[4] Because it carries equally these two oblations, the out-breathing and the in-breathing, the Samâna is he (the Hotri priest). The mind is the sacrificer, the Udâna is the reward of the sacrifice, and it leads the sacrificer every day (in deep sleep) to Brahman.

\para[5] There that god (the mind) enjoys in sleep greatness. What has been seen, he sees again; what has been heard, he hears again; what has been enjoyed in different countries and quarters, he enjoys again; what has been seen and not seen, heard and not heard, enjoyed and not enjoyed, he sees it all; he, being all, sees.

\para[6] And when he is overpowered by light, then that god sees no dreams, and at that time that happiness arises in his body.

\para[7] And, O friend, as birds go to a tree to roost, thus all this rests in the Highest Âtman, ---

\para[8] The earth and its subtile elements, the water and its subtile elements, the light and its subtile elements, the air and its subtile elements, the ether and its subtile elements; the eye and what can be seen, the ear and what can be heard, the nose and what can be smelled, the taste and what can be tasted, the skin and what can be touched, the voice and what can be spoken, the hands and what can be grasped, the feet and what can be walked, the mind and what can be perceived, intellect (buddhi) and what can be conceived, personality and what can be personified, thought and what can be thought, light and what can be lighted up, the Prâna and what is to be supported by it.

\para[9] For he it is who sees, hears, smells, tastes, perceives, conceives, acts, he whose essence is knowledge, the person, and he dwells in the highest, indestructible Self, ---

\para[10] He who knows that indestructible being, obtains (what is) the highest and indestructible, he without a shadow, without a body, without colour, bright, --- yes, O friend, he who knows it, becomes all-knowing, becomes all. On this there is this Sloka:

\para[11] He, O friend, who knows that indestructible being wherein the true knower, the vital spirits (prânas), together with all the powers (deva), and the elements rest, he, being all-knowing, has penetrated all."

\section*{FIFTH QUESTION}

\para[1] Then Saivya Satyakâma asked him: "Sir, if some one among men should meditate here until death on the syllable Om, what would he obtain by it?"

\para[2] He replied: "O Satyakâma, the syllable Om (AUM) is the highest and also the other Brahman; therefore he who knows it arrives by the same means at one of the two.

\para[3] If he meditate on one Mâtrâ (the A), then, being enlightened by that only, he arrives quickly at the earth. The Rik-verses lead him to the world of men, and being endowed there with penance, abstinence, and faith, he enjoys greatness.

\para[4] If he meditate with two Mâtrâs (A+U) he arrives at the Manas, and is led up by the Yagus-verses to the sky, to the Soma-world. Having enjoyed greatness in the Soma-world, he returns again.

\para[5] Again, he who meditates with this syllable AUM of three Mâtrâs, on the Highest Person, he comes to light and to the sun. And as a snake is freed from its skin, so is he freed from evil. He is led up by the Sâman-verses to the Brahma-world; and from him, full of life (Hiranyagarbha, the lord of the Satya-loka), he learns to see the all-pervading, the Highest Person. And there are these two Slokas:

\para[6] The three Mâtrâs (A+U+M), if employed separate, and only joined one to another, are mortal; but in acts, external, internal, or intermediate, if well performed, the sage trembles not.

\para[7] Through the Rik-verses he arrives at this world, through the Yagus-verses at the sky, through the Sâman-verses at that which the poets teach, --- he arrives at this by means of the Onkâra; the wise arrives at that which is at rest, free from decay, from death, from fear, --- the Highest."

\section*{SIXTH QUESTION}

\para[1] Then Sukesas Bhâradvâga asked him, saying: "Sir, Hiranyanâbha, the prince of Kosalâ, came to me and asked this question: Do you know the person of sixteen parts, O Bhâradvâga? I said to the prince: I do not know him; if I knew him, how should I not tell you? Surely, he who speaks what is untrue withers away to the very root; therefore I will not say what is untrue. Then he mounted his chariot and went away silently. Now I ask you, where is that person?"

\para[2] He replied: "Friend, that person is here within the body, he in whom these sixteen parts arise.

\para[3] He reflected: What is it by whose departure I shall depart, and by whose staying I shall stay?

\para[4] He sent forth (created) Prâna (spirit); from Prâna Sraddhâ (faith), ether, air, light, water, earth, sense, mind, food; from food came vigour, penance, hymns, sacrifice, the worlds, and in the worlds the name also.

\para[5] As these flowing rivers that go towards the ocean, when they have reached the ocean, sink into it, their name and form are broken, and people speak of the ocean only, exactly thus these sixteen parts of the spectator that go towards the person (purusha), when they have reached the person, sink into him, their name and form are broken, and people speak of the person only, and he becomes without parts and immortal. On this there is this verse:

\para[6] That person who is to be known, he in whom these parts rest, like spokes in the nave of a wheel, you know him, lest death should hurt you."

\para[7] Then he (Pippalâda) said to them: "So far do I know this Highest Brahman, there is nothing higher than it."

\para[8] And they praising him, said: "You, indeed, are our father, you who carry us from our ignorance to the other shore." Adoration to the highest Rishis! Adoration to the highest Rishis! Tat sat. Harih, Om!

\section*{Mundaka-Upanishad}

\section*{FIRST MUNDAKA}

\section*{FIRST Khanda}

\para[1] Brahmâ was the first of the Devas, the maker of the universe, the preserver of the world. He told the knowlede of Brahman, the foundation of all knowledge, to his eldest son Atharva.

\para[2] Whatever Brahmâ told Atharvan, that knowledge of Brahman Atharvan formerly told to Angir; he told it to Satyavâha Bhâradvâga, and Bhâradvâga told it in succession to Angiras.

\para[3] Saunaka, the great householder, approached Angiras respectfully and asked: "Sir, what is that through which, if it is known, everything else becomes known?"

\para[4] He said to him: "Two kinds of knowledge must be known, this is what all who know Brahman tell us, the higher and the lower knowledge."

\para[5] "The lower knowledge is the Rig-veda, Yagur-veda, Sâma-veda, Atharva-veda, Sikshâ (phonetics), Kalpa (ceremonial), Vyâkarana (grammar), Nirukta (etymology), Khandas (metre), Gyotisha (astronomy); but the higher knowledge is that by which the Indestructible (Brahman) is apprehended."

\para[6] That which cannot be seen, nor seized, which has no family and no caste, no eyes nor ears, no hands nor feet, the eternal, the omnipresent (all-pervading), infinitesimal, that which is imperishable, that it is which the wise regard as the source of all beings."

\para[7] As the spider sends forth and draws in its thread, as plants grow on the earth, as from every man hairs spring forth on the head and the body, thus does everything arise here from the Indestructible.

\para[8] The Brahman swells by means of brooding (penance); hence is produced matter (food); from matter breath, mind, the true, the worlds (seven), and from the works (performed by men in the worlds), the immortal (the eternal effects, rewards, and punishments of works).

\para[9] From him who perceives all and who knows all, whose brooding (penance) consists of knowledge, from him (the highest Brahman) is born that Brahman, name, form, and matter (food)."

\section*{SECOND Khanda}

\para[1] This is the truth: the sacrificial works which they (the poets) saw in the hymns (of the Veda) have been performed in many ways in the Tretâ age. Practise them diligently, ye lovers of truth, this is your path that leads to the world of good works!

\para[2] When the fire is lighted and the flame flickers, let a man offer his oblations between the two portions of melted butter, as an offering with faith.

\para[3] If a man's Agnihotra sacrifice is not followed by the new-moon and full-moon sacrifices, by the four-months' sacrifices, and by the harvest sacrifice, if it is unattended by guests, not offered at all, or without the Vaisvadeva ceremony, or not offered according to rule, then it destroys his seven worlds.

\para[4] Kâlî (black), Karâlî (terrific), Manogavâ (swift as thought), Sulohitâ (very red), Sudhûmravarnâ (purple), Sphulinginî (sparkling), and the brilliant Visvarûtpî (having all forms), all these playing about are called the seven tongues (of fire).

\para[5] If a man performs his sacred works when these flames are shining, and the oblations follow at the right time, then they lead him as sun-rays to where the one Lord of the Devas dwells.

\para[6] Come hither, come hither! the brilliant oblations say to him, and carry the sacrificer on the rays of the sun, while they utter pleasant speech and praise him, saying: "This is thy holy Brahma-world (Svarga), gained by thy good works."

\para[7] But frail, in truth, are those boats, the sacrifices, the eighteen, in which this lower ceremonial has been told. Fools who praise this as the highest good, are subject again and again to old age and death.

\para[8] Fools dwelling in darkness, wise in their own conceit, and puffed up with vain knowledge, go round and round staggering to and fro, like blind men led by the blind.

\para[9] Children, when they have long lived in ignorance, consider themselves happy. Because those who depend on their good works are, owing to their passions, improvident, they fall and become miserable when their life (in the world which they had gained by their good works) is finished.

\para[10] Considering sacrifice and good works as the best, these fools know no higher good, and having enjoyed (their reward) on the height of heaven, gained by good works, they enter again this world or a lower one.

\para[11] But those who practise penance and faith in the forest, tranquil, wise, and living on alms, depart free from passion through the sun to where that immortal Person dwells whose nature is imperishable.

\para[12] Let a Brâhmana, after he has examined all these worlds which are gained by works, acquire freedom from all desires. Nothing that is eternal (not made) can be gained by what is not eternal (made). Let him, in order to understand this, take fuel in his hand and approach a Guru who is learned and dwells entirely in Brahman.

\para[13] To that pupil who has approached him respectfully, whose thoughts are not troubled by any desires, and who has obtained perfect peace, the wise teacher truly told that knowledge of Brahman through which he knows the eternal and true Person.

\section*{SECOND MUNDAKA}

\section*{FIRST Khanda}

\para[1] This is the truth. As from a blazing fire sparks, being like unto fire, fly forth a thousandfold, thus are various beings brought forth from the Imperishable, my friend, and return thither also.

\para[2] That heavenly Person is without body, he is both without and within, not produced, without breath and without mind, pure, higher than the high Imperishable.

\para[3] From him (when entering on creation) is born breath, mind, and all organs of sense, ether, air, light, water, and the earth, the support of all.

\para[4] Fire (the sky) is his head, his eyes the sun and the moon, the quarters his ears, his speech the Vedas disclosed, the wind his breath, his heart the universe; from his feet came the earth; he is indeed the inner Self of all things.

\para[5] From him comes Agni (fire), the sun being the fuel; from the moon (Soma) comes rain (Parganya); from the earth herbs; and man gives seed unto the woman. Thus many beings are begotten from the Person (purusha).

\para[6] From him come the Rik, the Sâman, the Yagush, the Dîkshâ (initiatory rites), all sacrifices and offerings of animals, and the fees bestowed on priests, the year too, the sacrificer, and the worlds, in which the moon shines brightly and the sun.

\para[7] From him the many Devas too are begotten, the Sâdhyas (genii), men, cattle, birds, the up and down breathings, rice and corn (for sacrifices), penance, faith, truth, abstinence, and law.

\para[8] The seven senses (prâna) also spring from him, the seven lights (acts of sensation), the seven kinds of fuel (objects by which the senses are lighted), the seven sacrifices (results of sensation), these seven worlds (the places of the senses, the worlds determined by the senses) in which the senses move, which rest in the cave (of the heart), and are placed there seven and seven.

\para[9] Hence come the seas and all the mountains, from him flow the rivers of every kind; hence come all herbs and the juice through which the inner Self subsists with the elements.

\para[10] The Person is all this, sacrifice, penance, Brahman, the highest immortal; he who knows this hidden in the cave (of the heart), he, O friend, scatters the knot of ignorance here on earth.

\section*{SECOND Khanda}

\para[1] Manifest, near, moving in the cave (of the heart) is the great Being. In it everything is centred which ye know as moving, breathing, and blinking, as being and not-being, as adorable, as the best, that is beyond the understanding of creatures.

\para[2] That which is brilliant, smaller than small, that on which the worlds are founded and their inhabitants, that is the indestructible Brahman, that is the breath, speech, mind; that is the true, that is the immortal. That is to be hit. Hit it, O friend!

\para[3] Having taken the Upanishad as the bow, as the great weapon, let him place on it the arrow, sharpened by devotion! Then having drawn it with a thought directed to that which is, hit the mark, O friend, viz. that which is the Indestructible!

\para[4] Om is the bow, the Self is the arrow. Brahman is called its aim. It is to be hit by a man who is not thoughtless; and then, as the arrow (becomes one with the target), he will become one with Brahman.

\para[5] In him the heaven, the earth, and the sky are woven, the mind also with all the senses. Know him alone as the Self, and leave off other words! He is the bridge of the Immortal.

\para[6] He moves about becoming manifold within the heart where the arteries meet, like spokes fastened to the nave. Meditate on the Self as Om! Hail to you, that you may cross beyond (the sea of) darkness!

\para[7] He who understands all and who knows all, he to whom all this glory in the world belongs, the Self, is placed in the ether, in the heavenly city of Brahman (the heart). He assumes the nature of mind, and becomes the guide of the body of the senses. He subsists in food, in close proximity to the heart. The wise who understand this, behold the Immortal which shines forth full of bliss.

\para[8] The fetter of the heart is broken, all doubts are solved, all his works (and their effects) perish when He has been beheld who is high and low (cause and effect).

\para[9] In the highest golden sheath there is the Brahman without passions and without parts. That is pure, that is the light of light, that is it which they know who know the Self.

\para[10] The sun does not shine there, nor the moon and the stars, nor these lightnings, and much less this fire. When he shines, everything shines after him; by his light all this is lighted.

\para[11] That immortal Brahman is before, that Brahman is behind, that Brahman is right and left. It has gone forth below and above; Brahman alone is all this, it is the best.

\section*{THIRD MUNDAKA}

\section*{FIRST Khanda}

\para[1] Two birds, inseparable friends, cling to the same tree. One of them eats the sweet fruit, the other looks on without eating.

\para[2] On the same tree man sits grieving, immersed, bewildered by his own impotence (an-îsâ). But when he sees the other lord (îsa) contented and knows his glory, then his grief passes away.

\para[3] When the seer sees the brilliant maker and lord (of the world) as the Person who has his source in Brahman, then he is wise, and shaking off good and evil, he reaches the highest oneness, free from passions;

\para[4] For he is the Breath shining forth in all beings, and he who understands this becomes truly wise, not a talker only. He revels in the Self, he delights in the Self, and having performed his works (truthfulness, penance, meditation, \&c.) he rests, firmly established in Brahman, the best of those who know Brahman.

\para[5] By truthfulness, indeed, by penance, right knowledge, and abstinence must that Self be gained; the Self whom spotless anchorites gain is pure, and like a light within the body.

\para[6] The true prevails, not the untrue; by the true the path is laid out, the way of the gods (devayânah), on which the old sages, satisfied in their desires, proceed to where there is that highest place of the True One.

\para[7] That (true Brahman) shines forth grand, divine, inconceivable, smaller than small; it is far beyond what is far and yet near here, it is hidden in the cave (of the heart) among those who see it even here.

\para[8] He is not apprehended by the eye, nor by speech, nor by the other senses, not by penance or good works. When a man's nature has become purified by the serene light of knowledge, then he sees him, meditating on him as without parts.

\para[9] That subtle Self is to be known by thought (ketas) there where breath has entered fivefold; for every thought of men is interwoven with the senses, and when thought is purified, then the Self arises.

\para[10] Whatever state a man whose nature is purified imagines, and whatever desires he desires (for himself or for others), that state he conquers and those desires he obtains. Therefore let every man who desires happiness worship the man who knows the Self.

\section*{SECOND Khanda}

\para[1] He (the knower of the Self) knows that highest home of Brahman, in which all is contained and shines brightly. The wise who, without desiring happiness, worship that Person, transcend this seed, (they are not born again.)

\para[2] He who forms desires in his mind, is born again through his desires here and there. But to him whose desires are fulfilled and who is conscious of the true Self (within himself) all desires vanish, even here on earth.

\para[3] That Self cannot be gained by the Veda, nor by understanding, nor by much learning. He whom the Self chooses, by him the Self can be gained. The Self chooses him (his body) as his own.

\para[4] Nor is that Self to be gained by one who is destitute of strength, or without earnestness, or without right meditation. But if a wise man strives after it by those means (by strength, earnestness, and right meditation), then his Self enters the home of Brahman.

\para[5] When they have reached him (the Self), the sages become satisfied through knowledge, they are conscious of their Self, their passions have passed away, and they are tranquil. The wise, having reached Him who is omnipresent everywhere, devoted to the Self, enter into him wholly.

\para[6] Having well ascertained the object of the knowledge of the Vedânta, and having purified their nature by the Yoga of renunciation , all anchorites, enjoying the highest immortality, become free at the time of the great end (death) in the worlds of Brahmâ.

\para[7] Their fifteen parts enter into their elements, their Devas (the senses) into their (corresponding) Devas. Their deeds and their Self with all his knowledge become all one in the highest Imperishable.

\para[8] As the flowing rivers disappear in the sea, losing their name and their form, thus a wise man, freed from name and form, goes to the divine Person, who is greater than the great.

\para[9] He who knows that highest Brahman, becomes even Brahman. In his race no one is born ignorant of Brahman. He overcomes grief, he overcomes evil; free from the fetters of the heart, he becomes immortal.

\para[10] And this is declared by the following Rik-verse: "Let a man tell this science of Brahman to those only who have performed all (necessary) acts, who are versed in the Vedas, and firmly established in (the lower) Brahman, who themselves offer as an oblation the one Rishi (Agni), full of faith, and by whom the rite of (carrying fire on) the head has been performed, according to the rule (of the Âtharvanas)."

\para[11] The Rishi Angiras formerly told this true (science); a man who has not performed the (proper) rites, does not read it. Adoration to the highest Rishis! Adoration to the highest Rishis!
% !body-end
\end{document}
